% =========================================================
% CAPÍTULO 4: CONCLUSIONES
% =========================================================
\chapter{Conclusiones}
\label{cap:conclusiones}

% =========================================================
\section{Conclusiones generales}
\label{sec:conclusiones_generales}
% =========================================================

El objetivo de este trabajo fue dotar de control programático y
adquisición sincronizada a un arreglo fotoacústico preexistente del
Laboratorio de Fotofísica y Películas Delgadas del ICAT, cuya operación
manual exigía coordinar el disparo del láser, la adquisición del
osciloscopio y el registro de la temperatura de la muestra en cada
punto experimental. El sistema desarrollado cumple ese objetivo: los
tres instrumentos se operan desde una sola aplicación, la secuencia de
medición se ejecuta sin intervención entre puntos, y cada captura queda
guardada junto con las condiciones bajo las cuales fue tomada.

Conviene, sin embargo, precisar en qué consiste realmente la
aportación, porque la automatización de una tarea repetitiva no es por
sí sola un resultado interesante.

\subsection*{La trazabilidad como resultado principal}

Lo que la operación manual no garantizaba no era la velocidad, sino la
correspondencia entre un dato y las condiciones que lo produjeron. Una
captura guardada a mano queda descrita por lo que el operador alcanzó a
anotar, y esa anotación es independiente del archivo: puede perderse,
puede contradecirlo, y no hay forma de detectar la discrepancia después.

El sistema desarrollado resuelve eso escribiendo, en el mismo acto en
que guarda las muestras, los parámetros del láser vigentes, las escalas
consultadas al instrumento, la identificación del osciloscopio empleado,
la temperatura de la muestra y una marca explícita de si la medición se
tomó en condiciones anómalas. La consecuencia concreta de esa
decisión aparece en la \cref{sec:reconstruccion_escalas}: capturas
tomadas con bases de tiempo distintas convergen tras la reconstrucción
al mismo instante físico. Ese acuerdo no se obtiene por buena suerte,
sino porque los parámetros necesarios para reconstruir cada captura
viajan con ella.

La decisión de operar el osciloscopio en modo de solo lectura, expuesta
en la \cref{sec:decisiones}, es parte del mismo argumento. Un programa
que escribe la configuración del instrumento registra lo que creyó
haber establecido; uno que solo la consulta registra lo que el
instrumento realmente aplicó. La segunda opción entrega menos control y
más confianza, y para un sistema cuyo producto son datos destinados a
sostener afirmaciones, la confianza es lo que importa.

\subsection*{Sobre la validación sin instrumentos}

El acceso al laboratorio fue el recurso más escaso del proyecto. Esa
limitación obligó a una estrategia que resultó tener virtudes propias:
separar lo que depende de la lógica del programa, validable por
simulación y sesiones sintéticas, de lo que depende de la respuesta
efectiva del equipo.

La \cref{sec:validacion} documenta el alcance de esa separación.
Interesa destacar un punto metodológico. La herramienta de verificación
de sesiones no se dio por buena porque pasara sus pruebas, sino que se
sometió a una prueba de mutación en la que cada una de sus diez
comprobaciones se desactivó por separado para confirmar que su ausencia
hacía fallar el caso correspondiente. Esa prueba reveló, de paso, que el
código de salida del verificador podía permanecer en falla por una
condición ajena y dar por aprobado un caso con el chequeo apagado. Una
herramienta de auditoría que nunca se ha visto fallar no está
verificada; está sin poner a prueba.

\subsection*{Sobre las fallas documentadas}

El \cref{cap:resultados} dedica su sección más extensa a siete anomalías
observadas durante el desarrollo. Esa elección fue deliberada, y
conviene justificarla.

Seis de las siete no eran observables mientras ocurrían. Tres de ellas
tampoco lo eran inspeccionando los datos captura por captura: aparecieron
solo al contrastar el archivo con lo que su propio encabezado declaraba
contener, el código con su historial de cambios, y cada campo del archivo
con el lugar del código que lo produce. La más prolongada
—diez capturas consecutivas en las que no hubo adquisición— transcurrió
durante doce minutos frente a un operador atento, sin producir una sola
excepción.

De ahí se sigue el criterio con el que conviene juzgar un sistema de
este tipo. No es la ausencia de fallas, inalcanzable en un montaje que
integra tres instrumentos independientes con sus propios modos de
error. Es que las fallas resulten reconstruibles después. Todas las
anomalías descritas que se manifestaron pudieron caracterizarse meses
más tarde a partir de los archivos conservados, y esa posibilidad es en
sí misma un resultado del diseño: un sistema que hubiera guardado menos
las habría vuelto indetectables.

El caso de la propia herramienta de verificación, en la
\cref{sec:falla_umbral}, cierra el argumento por el lado incómodo. Un
criterio de umbral relativo produjo una diferencia aparente de treinta
microsegundos entre dos series de la misma sesión, que resultó ser un
artefacto del criterio y no un desplazamiento de la señal. El
instrumento de auditoría también se equivoca, y solo el contraste
contra una magnitud obtenida por una vía independiente permitió
distinguir el hallazgo del artefacto.

\subsection*{Limitaciones del trabajo}

Tres limitaciones deben quedar declaradas.

La primera es de alcance experimental: el modo de adquisición gobernado
por temperatura se ejercitó exhaustivamente por simulación y se
recorrió de punta a punta contra los tres instrumentos reales en una
sola sesión, sobre agua, en un intervalo de \SI{44}{\celsius} a
\SI{40}{\celsius} y bajo enfriamiento pasivo. La lógica está verificada
y el recorrido completo está demostrado; lo que no queda establecido es
su comportamiento sobre otros líquidos, intervalos más amplios o
regímenes de enfriamiento distintos del pasivo.

La segunda es que la calibración relativa entre los cuatro sensores de
temperatura quedó definida pero no aplicada, de modo que las lecturas
individuales arrastran la dispersión de fábrica del conjunto,
comparable en magnitud al umbral que gobierna el criterio de disparo.

La tercera es que el arreglo no mide la energía del pulso láser. En
consecuencia, las amplitudes registradas no son normalizables y el
sistema, por completo que sea su registro, no puede sostener por sí
solo afirmaciones cuantitativas sobre el coeficiente de absorción del
medio.

\subsection*{Balance}

El sistema entregado no caracteriza líquidos: hace posible
caracterizarlos. Convierte una sesión experimental que dependía de la
atención sostenida de un operador y de sus anotaciones en un
procedimiento que se ejecuta solo y que deja, en los propios archivos,
la evidencia necesaria para reconstruir después qué ocurrió, incluidas
las ocasiones en que algo salió mal. Las líneas de continuación
naturales de ese trabajo se enumeran a continuación.

% =========================================================
\section{Trabajo a futuro}
\label{sec:trabajo_futuro}
% =========================================================

El sistema desarrollado en este trabajo constituye una base funcional
que puede extenderse en varias direcciones. A continuación se describen
las líneas de trabajo identificadas durante el desarrollo.

\subsection*{La pregunta que los datos dejaron abierta}

El análisis de la sesión del 8 de septiembre de 2026, descrito en el
\cref{sec:falla_umbral}, dejó planteada una pregunta que este trabajo
no resuelve y que conviene enunciar con claridad, porque marca el
límite entre lo que la automatización establece y lo que no.

Descontado el artefacto del criterio de umbral, persiste un
desplazamiento del registro completo de aproximadamente
\SI{0.12}{\micro\second} a lo largo del barrido, en el sentido de
tiempos \emph{menores} conforme la muestra se enfría. La velocidad del
sonido en agua disminuye al bajar la temperatura en este intervalo, de
modo que el arribo debería retrasarse, no adelantarse; y la magnitud
esperada, para la separación hidrófono--haz registrada en los
metadatos de la sesión, es de unos \SI{0.08}{\micro\second}. El signo
es opuesto y el orden de magnitud no coincide.

Tres explicaciones quedan abiertas y este trabajo no dispone de
elementos para decidir entre ellas. La primera es instrumental: si la
fuente de disparo del osciloscopio dependiera de la propia señal
adquirida, el origen temporal se desplazaría con la amplitud y el
efecto sería enteramente artificial. La segunda es geométrica:
\SI{0.12}{\micro\second} equivalen a unas \SI{0.18}{\milli\meter} de
trayecto, una magnitud que la deriva térmica del arreglo o la ausencia
de una sujeción rígida del hidrófono pueden producir sin dificultad, y
que está muy por debajo de la incertidumbre con que esa distancia se
mide hoy. La tercera es de interpretación de la señal: el máximo se
ubica cerca de \SI{94}{\micro\second}, lejos del tiempo de arribo que
la separación registrada predice para la señal directa según el
\cref{ap:tiempos_arribo}, de modo que la característica que se está
siguiendo podría corresponder a una trayectoria reflejada y no a la
directa.

Distinguir entre las tres requiere intervenir el arreglo experimental
---fijar la geometría, verificar la fuente de disparo, caracterizar las
reflexiones--- y queda fuera del alcance de este trabajo, cuyo objeto
es la automatización del sistema y no la caracterización del fenómeno.
Se enuncia aquí porque su sola formulación depende de que el registro
conserve, para cada captura, las escalas del instrumento, la
temperatura de la muestra y la geometría de la sesión. La pregunta es
un producto del sistema, no una objeción a él.

\subsection*{Mejoras al hardware}

El módulo de temperatura actual emplea cuatro sensores DS18B20
conectados mediante buses 1-Wire independientes a un microcontrolador
ESP32 WROOM-32.
Una mejora natural sería sustituir la comunicación USB-serial por
una conexión inalámbrica (Wi-Fi o Bluetooth) que elimine el cable
entre el módulo y la computadora de control, simplificando el
arreglo experimental. Adicionalmente, el diseño de la PCB podría
incorporar un regulador de voltaje integrado que permita alimentar
el módulo directamente desde el puerto USB de la computadora de
control,
reduciendo el número de fuentes de alimentación independientes.

La calibración relativa entre los cuatro sensores de temperatura,
cuyo procedimiento se define en la \cref{sec:modulo_temperatura}, no
fue aplicada al firmware empleado en las mediciones reportadas. Su
ejecución es una tarea acotada ---una sesión de registro con los
cuatro sensores sumergidos en agua en agitación continua--- y su
resultado ajustaría las lecturas individuales al mismo orden de
magnitud que el umbral de decisión de la secuencia automática.

La reconexión automática del módulo de temperatura, documentada antes
como carencia, quedó implementada y verificada según la
\cref{sec:val_instrumento}. Subsiste un caso no cubierto: el reintento
opera sobre el puerto declarado en la configuración, de modo que una
reasignación del puerto por parte del sistema operativo sigue
requiriendo intervención manual. Resolverlo sin recorrer puertos ajenos
exigiría identificar al módulo por su descriptor de dispositivo antes
de abrir la comunicación.

En cuanto a la celda fotoacústica, trabajos futuros podrían explorar
geometrías optimizadas para minimizar reflexiones acústicas internas
y mejorar la relación señal a ruido de la señal detectada por el
hidrófono.

\subsection*{Mejoras al software}

Los metadatos de sesión incorporan la distancia entre el hidrófono y el
eje del haz y el volumen vertido, de modo que la ventana temporal útil
---que depende de ambos, según el \cref{ap:tiempos_arribo}--- puede
reconstruirse desde el archivo de datos y no desde la bitácora del
operador. Ambas magnitudes se capturan hoy a mano en la interfaz, y la
distancia se mide con flexómetro por no existir una sujeción fija del
hidrófono. Una montura que fije esa distancia, o su medición con un
instrumento de mayor resolución, reduciría la incertidumbre que hoy
domina el cálculo de la ventana.

El módulo de visualización actual muestra la forma de onda adquirida
tras concluir cada medición. Una extensión relevante sería implementar
un modo de previsualización en tiempo real durante la adquisición,
que permita al operador verificar la calidad de la señal antes de
confirmar el almacenamiento.

La prueba de mutación descrita en la \cref{sec:val_verificador} se
aplicó de forma manual y no existe en el repositorio como
procedimiento automatizado. Incorporarla como tal es lo que permitiría
que su resultado siga siendo válido cuando la herramienta de
verificación incorpore condiciones nuevas.

La biblioteca de control del láser mantiene su estado de conexión a
nivel de proceso y no por instancia. La aplicación no garantiza la
desconexión explícita al cerrarse, de modo que una terminación
irregular puede dejarla en un estado que impide la siguiente conexión
hasta reiniciar el proceso. Garantizar la desconexión en el cierre es
una corrección acotada, y conviene acompañarla de una comprobación de
conectividad independiente de la interfaz gráfica, como la que se
empleó para diagnosticar el caso.

La frecuencia de repetición del láser se emplea como constante del
programa para estimar el número de disparos integrados en cada
ventana, según se documenta en la \cref{sec:falla_laserhz}. El equipo
expone ese parámetro como registro consultable por la misma biblioteca
que el sistema ya utiliza. Leerlo en lugar de suponerlo, y guardar el
valor leído entre los metadatos de sesión, convierte un campo supuesto
en uno medido sin costo apreciable.

El módulo de almacenamiento guarda los datos crudos en archivos
\texttt{.npy} y los metadatos en un registro \texttt{.csv}.
Incorporar una base de datos ligera (por ejemplo, \texttt{SQLite})
facilitaría la búsqueda y comparación de mediciones por líquido,
concentración, temperatura o fecha sin necesidad de procesar los
archivos manualmente.

Finalmente, el pipeline de procesamiento podría integrarse
directamente en la aplicación: cálculo del coeficiente de absorción
óptica a partir de la amplitud de la señal fotoacústica, corrección
por la energía del pulso láser y generación automática de reportes
en formato PDF al concluir una sesión de medición.

\subsection*{Extensión a otros líquidos y concentraciones}

El sistema fue diseñado y validado con agua como líquido de
referencia. Una extensión directa consiste en caracterizar
soluciones acuosas de distintas concentraciones de absorbentes
de interés biológico o industrial, lo que permitiría construir
una base de datos de coeficientes de absorción en función de la
concentración y la temperatura.

La automatización del control de temperatura abre la posibilidad
de realizar barridos térmicos programados, obteniendo la dependencia
del coeficiente de absorción con la temperatura de forma sistemática
y reproducible, sin intervención manual entre puntos de medición.

\subsection*{Continuación académica}

Los resultados obtenidos en este trabajo sientan las bases para
una investigación de posgrado orientada a la caracterización
fotoacústica de mezclas de líquidos y al desarrollo de sensores
de concentración basados en el efecto fotoacústico pulsado.
La plataforma de automatización desarrollada puede reutilizarse
o adaptarse para otros experimentos del laboratorio que requieran
sincronización entre un láser pulsado y un sistema de adquisición
de señales.
